\section{The \texttt{did:lux:zoo}
 Method}

\subsection{DID Syntax}

Zoo DIDs follow W3C DID Core v1.0 \cite{w3c-did-core} with the method-specific identifier:

\begin{verbatim}
did:lux:zoo:<network>:<identifier>
\end{verbatim}

where \texttt{<network>} $\in \{\texttt{mainnet}, \texttt{testnet}, \texttt{devnet}\}$ and \texttt{<identifier>} is the Base58-encoded Blake2b-256 hash of the initial public key.

\begin{definition}[Zoo DID Document]
A Zoo DID Document $D$ is a tuple $(id, V, S, A, M)$ where:
\begin{itemize}
\item $id$ is the DID string,
\item $V = \{vk_1, \ldots, vk_m\}$ is the set of verification methods (public keys),
\item $S = \{s_1, \ldots, s_k\}$ is the set of service endpoints,
\item $A \subseteq V$ is the authentication key set,
\item $M$ is metadata including creation time, update history, and deactivation status.
\end{itemize}
\end{definition}

\subsection{CRUD Operations}

All DID operations are \IdentityVM{} transactions validated by consensus:

\begin{itemize}
\item \textbf{Create}: Initial document signed by genesis key. Valid if DID is unused, signature verifies, and document conforms to schema.
\item \textbf{Read}: Resolution queries the state trie with Merkle inclusion proofs for trustless light-client verification.
\item \textbf{Update}: Signed by a key in the document's \texttt{capabilityInvocation} set. Key rotation is atomic---old key authorizes new key in a single transaction.
\item \textbf{Deactivate}: Sets \texttt{deactivated} flag (irreversible). Document remains resolvable for credential verification but cannot authorize new operations.
\end{itemize}

